
This work introduces the δ(t) framework: a checkpoint linear probing protocol for measuring the temporal competition between spurious and core feature encoding during standard ERM training. Applied to Waterbirds/ResNet-50 over 30 epochs with three random seeds, the framework validates the existence of a statistically significant temporal gap (δ(t) > 0 for 13.3\% of training, p=0.022) driven by a feature complexity hierarchy (spurious 10× easier to linearly separate; all 3/3 complexity metrics directional, p<0.05 uncorrected) consistent with approximately 7× higher gradient signal for spurious features in early epochs (GDR=6.977). The transition epoch t\\textit{ has low cross-seed variance (std=2.0 epochs, CI upper bound=2.31 epochs), consistent with t\} being a structural property of SGD optimization on this task.

The hypothesis that DFR worst-group accuracy improvement would correlate positively with epochs trained past t\* is not supported (r=−0.815). DFR achieves WGA=0.806 at epoch 1 — before meaningful Waterbirds-specific training — and 0.871 at epoch 30. This suggests that ImageNet pretraining may contribute substantially to DFR robustness, though confirming this causal claim requires ablations not performed here.

These results are limited to a 30-epoch proof-of-concept on a single dataset and architecture. CelebA replication was not executed. Priority extensions are: (1) full 300-epoch training runs; (2) CelebA and text domain replication; (3) redesigned H-M4 using DFR absolute WGA; (4) annotation-free t\* detection via gradient norm proxies.

